\section{Related work}\label{sec:related}

\paragraph{Signal representations for EEG dementia classification.}
Spectral representations dominate: relative band power and EEG slowing are the most replicated AD
markers \cite{R3,F2,R15}. Time-domain and complexity features (Hjorth parameters, entropies) are
cheap and interpretable \cite{R15,R14}. Time--frequency transforms such as the continuous or
discrete wavelet transform retain non-stationary structure that a power spectrum averages away
\cite{R2,R16}, and several recent studies encode EEG as images for convolutional networks
\cite{R2,R6,R10}. Representation choice is thus treated as a first-class question --- but almost
always with convolutional models, and without varying the spatial structure the model sees.

\paragraph{Graph neural networks for EEG.}
Spectral graph convolution \cite{F13} and attention-based variants \cite{gat} underpin most EEG
GNNs, reviewed in \cite{F5}. For dementia, GNNs have been built on functional-connectivity graphs
\cite{F3}, learned adjacency \cite{F4}, multi-connectivity inputs \cite{R12}, time-resolved graphs
\cite{R10} and spatial--functional hybrids \cite{R9}; for Parkinson's disease, interpretable graph
learning has been used to localise abnormal connectivity \cite{F7}. The closest precedent \cite{F3}
compared connectivity measures as GNN inputs but held node features fixed. A 2025 preprint applied
a graph transformer to ds004504 for FTD versus controls with window-level prediction \cite{B6}.
None of these studies crosses node-feature representation with edge definition under a fixed model.

\paragraph{Graph construction for EEG.}
Edges may be defined spatially (inter-electrode distance), functionally (statistical coupling), by
learning, or by combinations of these. Phase-lag-based measures \cite{F11,wpli} suppress the
zero-lag coupling produced by volume conduction, which otherwise inflates functional graphs; weighted
graphs are recommended over binarised ones for EEG \cite{F9}, and phase-synchronisation networks
differ between dementia subtypes \cite{R13,R16}. Spatial graphs carry no subject-specific
information on their own and act as a smoothing prior; whether that prior helps is an empirical
question we test directly.

\paragraph{Evaluation protocol and leakage.}
Leakage --- information from test data reaching training --- is a general failure mode of
machine-learning-based science \cite{F8}, and has been shown to inflate prediction performance in
connectome-based neuroimaging models \cite{L5}. Saeb et al.\ showed that when the use case is
prediction for new individuals, record-wise cross-validation, which lets one person's records appear
in both training and test sets, greatly overestimates accuracy relative to subject-wise
cross-validation \cite{L1}. In EEG the problem takes a specific form: windows cut from one recording
are highly correlated, so splitting \emph{windows} rather than \emph{subjects} lets a model recognise
individuals instead of disease. Individual differences and non-stationarity make the partitioning
decision consequential for cross-participant EEG models, whose error rates improper partitioning can
understate severalfold \cite{L4}. Brookshire et al.\ found that segment-based holdout strongly
overestimated performance on unseen subjects in an Alzheimer's and an epilepsy dataset, and that most
translational deep-learning EEG studies they surveyed used it \cite{L2}; a recent audit across ADHD
and Parkinson's cohorts --- the latter the same cohort we use for external validation --- reproduced
the inflation in both \cite{L3}. In connectome-based neuroimaging models, leakage through repeated
subjects drastically inflated prediction performance, and more so in small datasets \cite{L5}. The
2026 scoping review of ds004504 attributes the upper end of reported accuracies on this benchmark to
the same effect \cite{R1}. Best-practice guidance for prediction studies \cite{L7} and the large
error bars of cross-validation in small samples \cite{L6} add a separate caution for cohorts of our
size. We measure the leakage effect directly on one pipeline in both cohorts, and report
subject-level confidence intervals throughout.

\paragraph{EEG and the blood--brain barrier.}
BBB breakdown precedes and predicts cognitive decline \cite{B5,B8}. Milikovsky et al.\ proposed
PSWEs as an EEG manifestation of BBB-related cortical dysfunction in AD and epilepsy \cite{B1}, with
a companion study tracing the mechanism to albumin-activated TGF$\beta$ signalling \cite{B2}.
Subsequent work found PSWEs uncommon in Parkinson's disease \cite{B3} and explored them as an
epilepsy biomarker \cite{B4}. Because ds004504 includes no BBB imaging, we treat PSWEs strictly as a
\emph{BBB-dysfunction-associated} EEG marker, not as a measurement of BBB integrity.

\paragraph{Positioning.}
Table~\ref{tab:positioning} summarises how this study differs from the closest prior work.

\begin{table}[t]
\centering\footnotesize
\caption{Positioning against the closest prior work on EEG dementia classification.
\checkmark\ = yes; -- = no; ? = not yet verified against the full text.}
\label{tab:positioning}
\begin{tabular}{lccccc}
\toprule
Study & Varies repr. & Varies graph & Fixed model & Subject-level split & Public data \\
\midrule
DICE-Net \cite{F2} & -- & -- & \checkmark & \checkmark & \checkmark \\
FC-measure GNN \cite{F3} & -- & \checkmark & \checkmark & ? & ? \\
Adaptive GCN \cite{F4} & -- & \checkmark & -- & ? & ? \\
TFR comparison \cite{R2} & \checkmark & -- & \checkmark & ? & \checkmark \\
Spatial--functional GNN \cite{R9} & -- & -- & -- & ? & ? \\
Multi-connectivity GCN \cite{R12} & -- & \checkmark & -- & ? & ? \\
Graph transformer \cite{B6} & -- & -- & \checkmark & \checkmark & \checkmark \\
\textbf{This study} & \checkmark & \checkmark & \checkmark & \checkmark & \checkmark \\
\bottomrule
\end{tabular}
\end{table}
% TODO (author): resolve every "?" from the full text before submission (several of these papers
% are paywalled; see paywalled_papers.md). Do not convert a "?" to a mark from the abstract alone.
